\documentclass[11pt]{amsart}
\usepackage[T1]{fontenc}
\usepackage{lmodern,amsmath,amssymb,mathtools,microtype,booktabs}
\usepackage[margin=1.15in]{geometry}
\usepackage[hidelinks]{hyperref}
\newtheorem{theorem}{Theorem}[section]
\newtheorem{proposition}[theorem]{Proposition}
\newtheorem{lemma}[theorem]{Lemma}
\newtheorem{corollary}[theorem]{Corollary}
\theoremstyle{definition}\newtheorem{definition}[theorem]{Definition}
\theoremstyle{remark}\newtheorem{remark}[theorem]{Remark}
\newcommand{\R}{\mathbb R}
\newcommand{\E}{\mathbb E}
\newcommand{\Prob}{\mathbb P}
\newcommand{\dd}{\,\mathrm d}
\newcommand{\one}{\mathbf 1}
\numberwithin{equation}{section}
\title[A2 v18: new proof extract]{Two-window reconstruction from intrinsic boundary laws}
\author{Qian Qi}
\date{September 28, 2026}
\begin{document}
\begin{abstract}
This extract contains the new proof chapter of A2 v18. It is not the
complete revised article or its preserved supplementary volumes.
Two active endpoint laws recover the stationary two-flight action;
its diagonal Hessian determines the two smooth boundary arcs.
A finite-bin experiment yields a fixed-exponent conditional inverse
for analytic periodic tables with quantitative asymmetry margins.
\end{abstract}
\maketitle
Endpoint grids are centered and invariant under simultaneous reversal
$(s,t)\mapsto(-s,-t)$, so that reversal permutes the finite cell outcomes.
\section{A two-window inverse for the full intrinsic law}
\label{sec:two-window}
The moment inverse and the full-law inverse answer different questions.
The former uses finitely many types of exact germs and makes explicit
which coefficient detects each contact jet. We now return to the
endpoint measures themselves. Two time windows suffice to recover their
generating action as a function, and a differential identity recovers
curvature without an infinite jet recursion. This gives smooth local
rigidity and an order-independent conditional stability estimate for
the analytic whole-table problem. Two density functions are, of course,
not two real observations. A finite experiment is given separately below.

\subsection{The intrinsic generating action}
Fix one orientation $b,o,b$ of an isolated normal channel. Write
$x(s)$ and $y(u)$ for unit-speed parametrizations of its source and
opposite boundary patches; both parameters increase in the common
transverse direction at the marked origins. Put
\[
 D(s,u)=|x(s)-y(u)|,\qquad
 \mathcal W(s,t)=D(s,u(s,t))+D(t,u(s,t)),
\]
where the intermediate point is selected by
$D_u(s,u)+D_u(t,u)=0$. At the normal orbit,
$D_{uu}=g^{-1}+\kappa_o>0$. The implicit-function theorem and channel
clearance therefore give a unique smooth physical branch on a rectangle
about $(0,0)$. The action is symmetric in $s,t$. All rectangles below
are reduced, when necessary, to stay in this branch.

\begin{lemma}[Affine dependence of the endpoint density on time]
\label{lem:tw-affine}
On a sufficiently small time collar and endpoint rectangle, the density
of the subprobability measure with respect to $\dd s\dd t$ is
\begin{equation}\label{eq:tw-density}
 f_T(s,t)=\frac{-\mathcal W_{st}(s,t)}{2\pi A}
                         (T-\mathcal W(s,t))_+.
\end{equation}
Here $-\mathcal W_{st}>0$. Suppose $T_2>T_1$ are in this collar and
$\mathcal W<T_1$ throughout a smaller rectangle $Q$. Then on $Q$,
\begin{align}
 \mathcal W&=\frac{T_1f_{T_2}-T_2f_{T_1}}{f_{T_2}-f_{T_1}},
                                                   \label{eq:tw-ratio}\\
 A&=\frac{(T_2-T_1)(-\mathcal W_{st})}
                         {2\pi(f_{T_2}-f_{T_1})}.  \label{eq:tw-area}
\end{align}
These expressions require neither a supplied area nor an exact onset.
\end{lemma}
\begin{proof}
In intrinsic section coordinates, the incoming momentum conjugate to
$s$ is $p=-\mathcal W_s$. Thus the section measure becomes
$(-\mathcal W_{st})\dd s\dd t$. The free residual time before the
first impact ranges from zero to $T-\mathcal W$. Taking the collar
shorter than the adjacent free flights excludes an earlier or a fourth
impact. Channel isolation places every relevant broken ray in the
chosen physical branch. Integrating residual time and dividing by
$2\pi A$ proves \eqref{eq:tw-density}. Positivity of the twist at the
normal orbit persists on a smaller rectangle. In the active rectangle,
write $w=-\mathcal W_{st}/(2\pi A)$. Then
$f_{T_j}=w(T_j-\mathcal W)$ and
$f_{T_2}-f_{T_1}=(T_2-T_1)w>0$. Elimination of $w$ gives
\eqref{eq:tw-ratio}, and recovery of $w$ gives \eqref{eq:tw-area}.
\end{proof}

The equality is between continuous density representatives. Equality of
the measures determines these representatives uniquely. Alternatively,
the support-onset function of the whole family recovers $\mathcal W$:
locally it is the infimum of those $T$ for which a point belongs to the
support. The two-window formula has the advantage that it differentiates
an ordinary quotient on a strictly active rectangle, not a moving
support boundary.

\begin{theorem}[Smooth local rigidity from two endpoint laws]
\label{thm:tw-smooth}
Two strictly active windows for one return orientation determine the
gap, the free area, and the two smooth boundary germs in their relative
placement, up to a Euclidean isometry. In particular, analyticity and
reflection symmetry are unnecessary for this local conclusion.
More explicitly, on the diagonal set
\begin{equation}\label{eq:tw-aux}
 \ell(s)=\tfrac12\mathcal W(s,s),\qquad
 a(s)=\mathcal W_s(s,s),\qquad v(s)=\sqrt{1-a(s)^2}.
\end{equation}
Near the normal origin $v>0$. The source curvature and the opposite
curvature at its nearest point are respectively
\begin{align}
 k(s)&=\frac{\mathcal W_{ss}(s,s)-\mathcal W_{st}(s,s)
                         -v(s)^2/\ell(s)}{v(s)}, \label{eq:tw-sourcecurv}\\
 k_o(u(s))&=\frac1{\ell(s)}
       \left(\frac{v(s)^2}{-2\ell(s)\mathcal W_{st}(s,s)}-1\right).
                                                    \label{eq:tw-othercurv}
\end{align}
Here $u(s)=u(s,s)$ is the opposite arclength parameter, with $u(0)=0$.
The gap is $g=\mathcal W(0,0)/2$, and the onset is $2g$.
\end{theorem}
\begin{proof}
On the diagonal the intermediate stationarity condition is $D_u=0$.
Strict convexity makes this the unique nearest point on the opposite
patch. Let $e=(x-y)/\ell$, and let $T=x'$ and $N$ be the source
unit tangent and inward normal, with $T'=kN$ and $N'=-kT$.
Near the normal orbit $e\cdot N>0$, so
\begin{equation}\label{eq:tw-frame}
                  e=aT+vN.
\end{equation}
Indeed the envelope identity gives $\mathcal W_s=D_s=e\cdot T=a$,
and $e$ has unit length. Choose the opposite tangent $T_o$ continuously
so that $T\cdot T_o=v>0$ near the marked origin. Since $e$ is the
opposite outward normal, direct differentiation of distance gives
\begin{equation}\label{eq:tw-distancehessian}
 D_{ss}=\frac{v^2}{\ell}+kv,\qquad
 D_{su}=-\frac v\ell,\qquad
 D_{uu}=\frac1\ell+k_o.
\end{equation}
For example, $D_{uu}=|T_o|^2/\ell-e\cdot y''$, and
$y''=-k_oe$ at the nearest point. These signs distinguish the
external dispersing geometry from an interior billiard.

The Hessian of the stationary two-flight action is the Schur complement
of the intermediate variable. On the diagonal it is therefore
\begin{equation}\label{eq:tw-schur}
 \mathcal W_{ss}=D_{ss}-\frac{D_{su}^2}{2D_{uu}},\qquad
 \mathcal W_{st}=-\frac{D_{su}^2}{2D_{uu}}
               =-\frac{v^2}{2\ell(1+\ell k_o)}.
\end{equation}
Subtract the first two expressions to obtain
\eqref{eq:tw-sourcecurv}; solve the last one for $k_o$ to obtain
\eqref{eq:tw-othercurv}.

For an explicit reconstruction put $x(0)=(0,0)$, $T(0)=(0,1)$ and
$N(0)=(-1,0)$. Solve the Frenet system with the recovered function $k$,
and integrate $x'=T$. This recovers the source arc uniquely. Equation
\eqref{eq:tw-frame} then recovers the opposite image as
\begin{equation}\label{eq:tw-foot}
             y(u(s))=x(s)-\ell(s)\{a(s)T(s)+v(s)N(s)\}.
\end{equation}
It is a genuine arc, rather than only a set with a possibly singular
parametrization: differentiating $D_u(s,u(s))=0$ gives
\begin{equation}\label{eq:tw-footspeed}
           u'(s)=\frac{v(s)}{1+\ell(s)k_o(u(s))}>0.
\end{equation}
In particular, $y(0)=(g,0)$ and the opposite tangent there is $(0,1)$.
Formula \eqref{eq:tw-area} recovers the area. The construction uses
functions on an interval, not only their Taylor series. Replacing
$(s,t)$ by $(-s,-t)$ reflects the entire reconstructed pair; there is
no independent reflection of its two members. All remaining choices
are a common translation and orthogonal transformation.
\end{proof}

At the normal origin, \eqref{eq:tw-schur} reduces to
$\mathcal W_{st}=-1/(2gc_o)$ and
$\mathcal W_{ss}-\mathcal W_{st}=c_b/g$. Thus the new reconstruction
agrees with the previously derived quadratic Hessian, including when
the two curvatures are equal. It also shows why the full endpoint law
contains more information than its two mass functions: the diagonal
mixed Hessian and the signed diagonal gradient are available before
any Taylor expansion is taken.

\subsection{Local stability without increasing the jet order}
Fix nested endpoint rectangles $Q_0\Subset Q$ containing the diagonal
segment under consideration. On the physical class in the next
proposition, all time, diameter and $C^3$ boundary bounds are uniform;
the gap, curvature, twist, and the length of the time separation have
positive lower bounds. Require also $T_1-\mathcal W\ge m>0$ on $Q$
and $v\ge v_*>0$ on its diagonal. These are open local geometric and
experimental margins, not knowledge of a boundary graph. The comparison
of arcs is in the canonical contact frames just used, and then in their
own intrinsic arclength coordinates on a fixed smaller interval.

\begin{proposition}[A Lipschitz inverse on physical density pairs]
\label{prop:tw-lipschitz}
Suppose two admissible physical pairs are observed at the same $T_1,T_2$,
with a common upper bound for the $C^2(Q)$ norms of their densities and
a common positive lower bound for $f_{T_2}-f_{T_1}$. Set
\[
 \delta=\max_{j=1,2}\|f_{T_j}-\widetilde f_{T_j}\|_{C^2(Q)}.
\]
After the permitted common reflection, the gap and area errors and
the $C^2$ errors of both intrinsic boundary arcs on a smaller fixed
interval are at most $C\delta$. The constant depends on the stated
margins and bounds, but not on a jet order.
\end{proposition}
\begin{proof}
Multiplication and reciprocal are locally Lipschitz on bounded subsets
of $C^2(Q)$ with denominator bounded below. Formula
\eqref{eq:tw-ratio} therefore bounds the $C^2$ action error by
$C\delta$. Equations \eqref{eq:tw-sourcecurv} and
\eqref{eq:tw-othercurv}, whose denominators have positive margins,
bound the source curvature error and the opposite curvature error
as functions of $s$ by $C\delta$. The gap and area follow from
\eqref{eq:tw-area} and the value at $(0,0)$.

Continuous dependence for the Frenet system, or its integral equation
and Gronwall's inequality, gives a $C^2$ error $C\delta$ for the source
arc in its intrinsic parameter. For the opposite arc one must not
differentiate \eqref{eq:tw-foot} twice and silently ask for a third
derivative of $\mathcal W$. Instead integrate
\eqref{eq:tw-footspeed}. Its speed is bounded above and below and
has error $C\delta$, so both $u(s)$ and its inverse have uniform
error $C\delta$ on smaller intervals. The physical $C^3$ boundary
bound makes opposite curvature uniformly Lipschitz in its arclength.
Consequently the two recovered opposite curvature functions, now
compared at the same $u$, differ by $C\delta$. Apply the Frenet
system in this intrinsic parameter with the common initial tangent
and the recovered gap. This gives the asserted $C^2$ opposite-arc
estimate using only a $C^2$ action error. No differentiability of
an arbitrary noisy curvature estimate is assumed.
\end{proof}

\subsection{Finite bins and a fixed-exponent whole-table bound}
We spell out the sampling step, since a norm on density functions is
not itself a finite observation. Let $0<\beta\le1$. Assume the
admissible active densities have a uniform $C^{2,\beta}(Q)$ bound.
Divide a slightly larger fixed rectangle into squares of side $h$.
For each cell $B$ observe its unconditional probability
$\mu_{b,T_j}(B)$; keep failure and success outside the rectangle as
additional outcomes. Thus all preparations are charged. The collection
of these probabilities is a finite multinomial observation of the same
endpoint-arclength law, not a new trajectory sensor.

\begin{lemma}[Recovering two derivatives from cell probabilities]
\label{lem:tw-bins}
There are linear reconstructions from finitely many cell probabilities
such that, if each probability has error at most $\epsilon$, their
reconstructed density $\widehat f$ satisfies
\begin{equation}\label{eq:tw-bins}
 \|\widehat f-f\|_{C^2(Q_0)}
             \le C\{h^\beta+\epsilon h^{-4}\}.
\end{equation}
In particular, if two functions in the prior class have cell
probabilities differing by at most $\epsilon$, their $C^2(Q_0)$
distance satisfies the same estimate, with a changed constant.
There are $O(h^{-2})$ observed cells per time window.
\end{lemma}
\begin{proof}
Use normalized coordinates at a grid cell. The averages of
$1,z,z^2$ on the three unit intervals centered at $-1,0,1$ form the
invertible matrix with rows $(1,i,i^2+1/12)$, $i=-1,0,1$.
Its tensor square therefore recovers a polynomial of coordinate degree
at most two from a $3\times3$ block of cell averages. The corresponding
linear operator and its derivatives through order two are bounded
on each fixed normalized block.

For a $C^{2,\beta}$ function its degree-two Taylor polynomial at the
block center has remainder $O(h^{2+\beta})$. The recovered local
polynomial reproduces that Taylor polynomial exactly. Its derivative
errors of order $j\le2$ are consequently $O(h^{2+\beta-j})$.
A probability error $\epsilon$ becomes an average-density error
$\epsilon h^{-2}$; differentiating the recovered polynomial $j$
times multiplies its error by $O(h^{-j})$. Patch these local
polynomials with a smooth locally finite partition of unity subordinate
to fixed enlargements of the cells. Its derivatives of order $j$ are
$O(h^{-j})$, and the overlap number is uniformly bounded. Applying
Leibniz's rule to the local polynomial minus $f$ gives
$O(h^{2+\beta-j}+\epsilon h^{-2-j})$ for every $j\le2$.
The larger observation rectangle avoids boundary stencils. Taking
$j=2$ proves \eqref{eq:tw-bins}. Applying the construction to both
functions and subtracting proves the comparison assertion.
\end{proof}

For the next theorem fix a connected finite channel network visiting
all labelled obstacles, integer edge copy labels, and two independent
fundamental cycle labels. These are discrete markings; signed
cross-channel registrations and relative edge signs are not supplied.
On each obstacle require a uniformly bounded holomorphic support
function on a strip of positive width, positive curvature and geometric
clearance. After Steiner centering write the support function as $p$.
For a fixed $\eta>0$ require
\begin{align}
 \|p(\cdot+\phi)-p\|_2&\ge
             \eta\operatorname{dist}(\phi,2\pi\mathbb Z),
                                                   \label{eq:tw-rot}\\
 \|p(-\cdot+\phi)-p\|_2&\ge\eta.    \label{eq:tw-ref}
\end{align}
Include bounded diameters, bounded lattice and inverse, positive area,
the uniform local margins above, and a known bounded bracket for each
selected-channel onset. Take a compact physical class with these
bounds in the $C^2$ table metric. The uniform analytic bounds imply the
fixed $C^{2,\beta}$ density bounds after the collars are reduced.
The table metric is the maximum labelled-support $C^2$ error and marked
lattice error, minimized over one common Euclidean isometry, with tree
representative-copy labels gauged to zero.

\begin{theorem}[A fixed-exponent finite experiment for the whole table]
\label{thm:tw-global}
On this class there are $C>0$, $0<\theta<1$ and $\epsilon_0>0$
with the following property. A finite pilot followed by two active
endpoint histograms per edge gives an admissible reconstruction with
\begin{equation}\label{eq:tw-global}
 \operatorname{dist}_{\rm tab}(\widehat{\mathcal T},\mathcal T)
 \le C\{h^\beta+\epsilon h^{-4}\}^{\theta},
\end{equation}
provided every reported cell probability and pilot mean has error at
most $\epsilon<\epsilon_0$. The pilot and histogram windows use
absolute times. Neither the exact onsets, leading geometry, obstacle
images nor contact registrations are supplied. Every edge has its own
simultaneous reversal orbit, fixed across that edge's two windows.

With $h=\epsilon^{1/(\beta+4)}$ this gives
\begin{equation}\label{eq:tw-rate}
 \operatorname{dist}_{\rm tab}(\widehat{\mathcal T},\mathcal T)
                  \le C\epsilon^{\beta\theta/(\beta+4)}.
\end{equation}
For independent fresh phase preparations, simultaneous accuracy holds
with probability at least $1-\delta$ using at most
\begin{equation}\label{eq:tw-budget}
 C|E|\epsilon^{-2}
       \log\left(\frac{C|E|(1+h^{-2})}{\delta}\right)
\end{equation}
preparations, including failures. The constants in this theorem are
independent of a contact-jet order. The bound is an upper bound for
this finite marked experiment; no minimax or count-only rate is claimed.
\end{theorem}
\begin{proof}
\emph{Pilot and active rectangles.}
Uniform local geometry gives a collar of length $d_*>0$ and constants
$c,C>0$ such that $cd^2\le P_b(T_0+d)\le Cd^2$ for
$0<d<d_*$. Before $T_0$ the selected event has zero mass. Choose
once and for all a positive threshold and grid spacing sufficiently
small compared with these constants and $d_*$. Query the known onset
bracket, adding a fixed short collar after its upper endpoint. The
first reported mass above the threshold occurs at a time $\tau$ with
\[
                         0<\tau-T_0<d_*/16
\]
when the mean errors are sufficiently small. To be explicit, choose
$\lambda>0$ with $\sqrt{4\lambda/c}<d_*/32$, use grid spacing
less than $d_*/32$, and trigger at reported mass $2\lambda$ with
errors below $\lambda/4$. A grid point at offset at least
$\sqrt{4\lambda/c}$ and less than that number plus the grid spacing
must trigger, whereas a pre-onset point cannot. There are only a fixed
number of pilot queries, depending on the prior class, not on
$\epsilon$ or $h$.

Use $T_1=\tau+d_*/4$ and $T_2=\tau+d_*/2$. Choose $Q$ so small
that $\mathcal W-T_0\le d_*/8$ on it for every admissible channel.
Then both times stay inside the common collar, and
$T_1-\mathcal W\ge d_*/8$. The time separation and twist have
positive margins. The same argument holds for every physical candidate
compatible with the pilot data, with slightly enlarged error constants.
No global monotonicity of the three-impact probability is used.

\emph{Local and analytic reconstruction.}
Apply Lemma~\ref{lem:tw-bins} and Proposition~\ref{prop:tw-lipschitz}
to two compatible physical candidates. Their entire contact-arc pairs
are $O(e)$ apart in intrinsic $C^2$ coordinates, where
$e=h^\beta+\epsilon h^{-4}$. Positive curvature converts each arc
to a common interval of normal angles, with an $O(e)$ support-function
error. Indeed the angle map has derivative equal to curvature, bounded
away from zero; its inverse and the corresponding support values and
their first two derivatives depend locally Lipschitzly on the arc,
using the uniform third-derivative bound.

For clarity the continuation step does not invoke convergence of a jet
recursion. A difference of holomorphic support functions is bounded on
a fixed periodic strip and is $O(e)$ on a fixed real subinterval. Slit
a narrower strip cylinder along a shorter such interval. The harmonic
measure of the slit has a positive minimum $\theta$ on a further
narrower closed cylinder, including the slit by continuity. The
subharmonic maximum principle applied to the logarithm of the modulus
gives an $O(e^\theta)$ bound there. Cauchy estimates on fixed disks
give the same bound through two derivatives on the real circle.
The slit endpoints are regular boundary points, and zeros are handled
by subharmonic regularization. Thus each entire analytic pair image
is $O(e^\theta)$-close in its own contact frame.

\emph{Unregistered assembly.}
Center every recovered obstacle at its Steiner point. Two approximate
matches to a placed copy differ by an orthogonal transformation whose
action on its support function is $O(e^\theta)$. Inequality
\eqref{eq:tw-ref} excludes the incorrect reflection component once
the error is small. Inequality \eqref{eq:tw-rot} bounds the relative
rotation angle by $O(e^\theta/\eta)$. Translations are controlled
by the Steiner points. Iterating along the fixed spanning tree places
all representative obstacles with error $O(e^\theta)$, without any
supplied relative sign or contact offset. For a non-tree edge the
translated target image determines a displacement $d_e=Lk'_e$.
For two independent cycle labels put
$\mathsf K=[k'_{e_1}\ k'_{e_2}]$. Then
\[
 L=[d_{e_1}\ d_{e_2}]\mathsf K^{-1},\qquad
 \|L-\widetilde L\|\le C\|\mathsf K^{-1}\|e^\theta.
\]
Remaining cycles and physical constraints only restrict the candidates.
This proves \eqref{eq:tw-global}; a bounded prior diameter handles
errors above the small-error regime. The minimizing comparison uses
one global isometry, not independent final placements.

\emph{An admissible estimator.}
Retain the pilot observations and all histogram probabilities. Minimize
the maximum discrepancy over the compact physical class, allowing one
common reversal per edge across both histograms. A first-eligible-point
selection from a fixed countable dense family, within $\epsilon$
of the infimum, is Borel. The probability of a cell and each pilot
mean is continuous on the physical class: grazing, patch-boundary,
cell-boundary and terminal-time events have phase volume zero, and
outside these events finite collision records vary continuously.
Uniform separation bounds their length on the queried time interval,
so dominated convergence applies. The true table is compatible on the
mean-accuracy event; the selected table consequently has discrepancy
at most a fixed multiple of $\epsilon$. The preceding comparison
therefore applies to it. This is a regularization theorem, not a claim
of an efficient search over analytic tables.

\emph{Preparation count.}
One phase preparation at a histogram time produces one multinomial
outcome, hence simultaneously one Bernoulli indicator for every cell.
Hoeffding's inequality and a union bound over
$O(|E|(1+h^{-2}))$ cells and pilot channels give accuracy $\epsilon$
with $C\epsilon^{-2}\log(C|E|(1+h^{-2})/\delta)$ independent
preparations at each queried time. The pilot has only a prior-dependent
number of times per edge. Conditional on its record, the two histogram
times are fixed; using fresh preparations justifies the same bound for
this adaptive stage. There is no additional factor $h^{-2}$ in the
preparation count, since the same sample records all cell indicators.
This proves \eqref{eq:tw-budget}. Balancing the two terms in $e$
proves \eqref{eq:tw-rate}.
\end{proof}

The new theorem uses the original endpoint law, not a marked length
spectrum, and it does not remove the discrete lattice-copy marking.
It removes the need for an all-order stability estimate when the full
law is retained. The moment-compressed inverse, its dimension-sharp
finite-window results, the registered symmetric-table theorem, and
the smooth relative-law theory remain separate retained conclusions.
Smooth local rigidity here is an equality-of-functions statement; it
has no implication that a formal count fiber converges analytically,
or that equality of all smooth count jets is equality of count germs.

\end{document}
